% ============================================================
% DEFERRED REVISION — LAB-KEYSTONE DEPENDENCY (ruled 2026-07-29)
%   STATUS: queued, NOT to be executed until the laboratory boredom
%   section is drafted. Author ruling: working through the lab material
%   will likely revise the hypotheses and interpretations, so these
%   spans are revised ONCE, afterward, against settled findings —
%   not now, and not twice.
%
%   WHAT IS STALE: the rendered prose below cites the laboratory
%   keystone in its SUPERSEDED form — a brain measure reading
%   boring-like against a gaze measure reading engaged. The O-9
%   reprocessing (2026-07-16, N=8) withdrew it: P1's "clinical EEG
%   ~ boring" came from one clean subject and does not replicate
%   (clinical DMN is LOW = engaging-like; boring-like in 4/8 = chance).
%   The corrected keystone is ALL-PHYSIOLOGY-ENGAGED vs
%   SELF-REPORT-BORED (clinical self-report 6.0/9).
%
%   LAYER ASSIGNMENT (author-ruled 2026-07-29, governs the rewording):
%   physiology as a whole = the OCCUPIED SURFACE (outward comportment);
%   the post-video retrospective self-report = the HOLLOW DEPTH.
%   Grounded in fcm-05 SS24-28: at the dinner party we are "totally
%   involved" and DENY boredom, recognizing it only retrospectively;
%   "behavioural engagement is not evidence against boredom; in this
%   form it is the very medium of it."
%
%   SITES IN THIS FILE: see the per-site marker(s) below.
%   COMPANION SITES: DESKTOP-M3-DRAFT-v1.tex (M3.6) and
%   DESKTOP-M4-DRAFT-v1.tex (M4.1 P6); DESKTOP-SECTION-OVERLEAF-v1.tex
%   mirrors BOTH and is GENERATED — re-flatten after the edits, then
%   recompile DESKTOP-SECTION-ASSEMBLED-v1.tex (it \input-s M3/M4, so
%   the 37-pp PDF carries the stale claim until then).
%
%   CANONICAL SOURCES WHEN THE TIME COMES:
%     corpus/index/Heidegger - The Fundamental Concepts of Metaphysics/
%       _synthesis/fcm-king-salvo-bridge.md  (S4, revised 2026-07-29)
%     tmp/Dissertation/Part_III/reanalysis/boredom-o9-physiological-findings.md
%   NOTE: the structural point at both sites SURVIVES unchanged (a
%   survey holds one channel; the laboratory holds two). Only the
%   naming of the two channels is wrong.
% ============================================================
% ============================================================
% Part III — Desktop Section — M4 (Phenomenological Analysis: The Stalled
%   Chain and the Answered Question) — DRAFT v1
% Mode: Stage 3 guided drafting (per-paragraph protocol). M4.1 architecture
%   (7 ¶¶, 3 batches) approved 2026-07-27; batch 1 = ¶¶1–3 DRAFTED
%   2026-07-27 — IN REVIEW. Band G throughout; style canon = approved M3.1
%   sample (humanistic, slowed statistics, one claim per sentence, quotes
%   worked one at a time). Conventions as DESKTOP-M3-DRAFT-v1.tex (MLA;
%   respondent-code parentheticals; chain nodes A\textsubscript{n}; no
%   module labels in rendered text; no metatextual forward-refs; per-student
%   percentages w/ % symbol; no price/cost/ledger metaphors; student typos
%   verbatim, no [sic]; "YouTube platform/route"; terminology lock
%   [medium-supply vs phantasia gap-filling only]; "being-in" hyphenated).
%
% SOLICITATION RULING (2026-07-27, user-confirmed — gate CLOSED; full text
%   in HANDOFF-M4-GATE-PRESENTED-2026-07-25.md §5): solicitation in the
%   course deployment splits by agent and moment. (1) The recruiting hail
%   is INSTITUTIONAL and pre-perceptual — the professor's stated pitch at
%   the course's outset, in language, before the world renders; its content
%   = the VLE-only affordance cluster (co-watching the assigned procedure's
%   footage in shared virtual space + communication, toward group analysis).
%   (2) The world's own solicitation work is SUSTAINING — holding attention,
%   charging the clinical image, keeping use alive under mandate — and goes
%   silent only where soliciting would have to become TASKING (A3→A4).
%   Part II untouched (A1 = the Sonny event's own stationing: predatory-
%   solicitation mimicry atop the generic leisure entry-motive). The
%   outline's "solicited at A2/A3" = compression, unfolded in ¶¶1–2.
%   Leergelassenheit gloss CORRECTED: the TASK is withheld, not
%   solicitation; the emptiness is REGIONAL, stationed at the last link.
%   Hingehaltenheit sharpened: mandated duration = entry-under-mandate-
%   without-leisure-motive from inside. Drafting precision: the pitch
%   promised a place FOR the (course-side) task, not a task inside the
%   world — why the silence could hide under a successful pitch.
%
% AUTHOR THEORY CAUTION (2026-07-27): boredom is ROUTE-INDIFFERENT at the
%   content level — footage that bores, bores on YouTube and in the VLE
%   alike; the routes affect how fast and to what degree, not whether.
%   Boredom therefore cannot ground the route-differential (that work
%   belongs to the WE, M4.3); its diagnostic value here = the watching-only
%   situation both routes stage. The WE-reduces-boredom effect is
%   SPECULATIVE ONLY (no reliable data) — if voiced at all (M4.3/M5
%   territory), it must be marked as speculation. User 2026-07-27: full
%   three-forms exposition PROCEEDS as planned (¶¶4–5); removable or
%   editable at review.
%
% INHERITANCE (build on, never re-derive): M3.3 ¶¶4–5 (A3 route-selection;
%   seven-causes individuation) · M3.4 (incorporation-debt; matched-numbers
%   insight) · M3.5 (floors 1/11/4/13, awareness-gap 12; coordination-
%   demand theorizing deferred to M4.3) · M3.6 (promissory footnote: the
%   three-forms exposition due in this section — to be paid at ¶¶4–5).
%   ¶2's "three of every five" inherits M3.2's D2 per-student stat (147/241
%   = 61%); presence-from-those-who-entered inherits F7. W26-R037 fragments
%   re-quoted deliberately for continuity (standing preference: make
%   same-student threads loud), char-exact verified 2026-07-24 per the
%   M3.6 ledger ("then" typo verbatim).
%
% DECISION POINTS RESOLVED 2026-07-27: pitch = author-confirmed course
%   fact, first-person plural ("we presented"); published trace footnoted =
%   ASEE 2025 "Scalable Virtual Reality" (King, Diec & Salvo; cite by PDF
%   page — conference paper, no journal pagination; vle-05 L1). VERIFIED
%   this session: span archon exact-1.00+bbox (vle-05 ledger) AND pdftotext
%   page-pinned to PDF p. 4 ("supported real-time voice chat and screen
%   sharing, enabling collaborative team viewing if students wanted to meet
%   in groups within the virtual environment"). OPEN AT REVIEW: the pitch's
%   delivery vehicle (syllabus vs first-session announcement; whether each
%   term) not yet author-specified — ¶1 kept to "from the outset"; confirm
%   or tighten wording. EXPANSION CANDIDATE held in reserve: vle-05 L4 —
%   the authors' own published boredom line ("engagement alerts will be
%   integrated to address student feedback on boredom," PDF pp. 9–10) —
%   fits ¶2 or ¶7 if wanted. WORKS CITED ADD QUEUED: King, Christine E.,
%   Kadin Diec, and Dalton Salvo. "Scalable Virtual Reality Global Clinical
%   Immersion for Culturally Responsive Engineering Design Skills." ASEE
%   Annual Conference & Exposition, Paper ID #49715, ASEE, 2025.
% Batch-2 verification ledger (2026-07-27, index-entry-first then PDF —
%   corpus/metaphysics/Heidegger…Fundamental Concepts…pdf; OCR text layer
%   w/ minor artifacts [stray pipes]; entry pdf-page loci run 1–2 HIGH in
%   the early chapters — true pdf pages pinned below; BOOK pages cited in
%   prose per the entry's triple loci):
%   "that which holds us in limbo and yet leaves us empty" p.87 (pdf 55) ✓
%   "the tasteless station of some lonely minor railway" + "It is four
%     hours until the next train arrives" + "catch ourselves looking at
%     our watch yet again" p.93 (pdf 58) ✓
%   "helpless gesture" p.97 (pdf 60) ✓
%   "There is nothing at all to be found that might have been boring
%     [about this evening]" + "I was bored after all this evening" p.110
%     (pdf 66) ✓
%   "what is boring does not come from outside: it arises from out of
%     Dasein itself" p.128 (pdf 76) ✓
%   "it is boring for one" + "'it is boring for one' to walk through the
%     streets of a large city on a Sunday afternoon" p.135 (pdf 79) ✓
%   "compelled to listen" p.136 (pdf 80) ✓
%   "the while becomes long" p.152 (pdf 88) ✓
%   Leergelassenheit / Hingehaltenheit / Zeitvertreib ✓ (translators'
%     brackets + glossary).
%   GERMAN FORM-NAMES (Gelangweiltwerden von etwas / Sichlangweilen bei
%     etwas / es ist einem langweilig): NOT PRINTED in the McNeill–Walker
%     translation (GA-side apparatus only) → per M3.6 precedent the prose
%     uses ENGLISH form-names + the translation's own "it is boring for
%     one." WORKS CITED ADD QUEUED: Heidegger, The Fundamental Concepts
%     of Metaphysics: World, Finitude, Solitude, trans. McNeill & Walker,
%     Indiana UP, 1995.
% Batch-2 v2 RULINGS APPLIED (user review 2026-07-27, "P4/P5"):
%   (1) opening reworded per user: "one of the most sustained,
%       philosophical, and structural analyses it has received";
%   (2)–(5) + P5: exposition RESTRUCTURED ¶¶4–5 → ¶¶4–7, one ¶ per form,
%       each with direct definition + Heidegger scene EXPLAINED (not
%       merely quoted) + demonstrative example + evidence requirement
%       worked against the record. "Channel" defined plainly up front
%       (¶4). "Named culprits/frictions/exits" unpacked as the form's
%       anatomy — culprits (quality/length/lag), occasions (crash/
%       buffering/wayfinding), counter-move = the flight to the same
%       footage elsewhere — w/ verified quotes. Second form: structure
%       spelled out + streaming-evening example + lab two-channel
%       demonstration (reference only) + record's echo-shape w/ codes.
%       Third form: Sunday-afternoon scene EXPLAINED (works because
%       nothing in it is boring) + student-version example + "powerless"/
%       "compelled to listen" + stillness/felt-time evidence + no-instance
%       disclaimer. MODULE NOW 9 ¶¶ (¶¶8–9 = application + diagnosis).
%   Batch-2 v2 additional verifications (2026-07-27): "all passing the
%   time is powerless against this boredom" p.135 (pdf 79) ✓ ·
%   W25-R017/Q9 "the youtube format ended up being simpler to use and
%   watch" char-exact vs open_text.csv ("youtube" lowercase verbatim) ✓ ·
%   W25-R014/Q9 re-confirmed char-exact vs open_text.csv ("dont" typo
%   verbatim) ✓.
% ============================================================
% --- M4.1 Completion, not bandwidth (¶¶1–7 APPROVED 2026-07-27 (¶¶4–7
%     approved w/ two ordered changes, both APPLIED same day): (a) RULE
%     RE-AFFIRMED — never open a sentence with "And" (three violations
%     fixed: ¶3 [approved-text touch, wording only], ¶4, ¶5); (b) P6
%     ruling — W26-R016/Q6 + S25-R009/Q10 quoted IN FOOTNOTE at the
%     culprit-less-responses clause w/ insight explanation (verbatims per
%     2026-07-24 ledger). ¶4 = apparatus + channel term; ¶5 = first form;
%     ¶6 = second form; ¶7 = third form (FCM §§19–23/24–28/29–36;
%     vocabulary from the FCM entry + crosswalk §1/§4, NOT Part II v4;
%     English form-names per verification finding). ¶¶8–9 APPROVED
%     2026-07-27 — **M4.1 COMPLETE (9 ¶¶, all approved)**. ¶8 = the two
%     moments at the deployment's
%     address (Leergelassenheit/Hingehaltenheit named; at-hand emptiness;
%     station-refusal analogy → hospital-refuses-itself-as-clinic;
%     S25-R035/Q6 orekton-withheld continuity quote; mandate-held close);
%     ¶9 = mis-attunement diagnosis (learner/ecology/task) + energeia
%     atelēs + regional emptiness + route-indifference guard + hinge to
%     the for-the-sake-of-which (M4.2's demand, unnamed). MODULE FULLY
%     DRAFTED pending ¶¶8–9 review.
%     Batch-3 verification ledger (2026-07-27): "To leave empty means to
%     be something at hand that offers nothing" p.103 (pdf 63) ✓ · "leave
%     us empty precisely because they are at hand" p.102 (pdf 63) ✓ ·
%     "The station at hand refuses itself to us as a station and leaves
%     us empty because the train that belongs to it has not yet arrived"
%     p.103 (pdf 63) ✓ · S25-R035/Q6 char-exact vs open_text.csv ✓
%     (continuity re-quote; also in M3.2 ledger) · energeia atelēs =
%     Part I canonical (spelling per §1.5 v3; gloss "incomplete
%     actuality, actuality whose end lies beyond it" reused verbatim from
%     §1.1 v3; locus Physics III.1, 201a10 per §1.1 fn; Metaph. IX.6
%     1048b30–36 available via §1.5). PHILOSOPHICAL GUARD: energeia
%     atelēs attached to interest-underway-toward-doing, NEVER to
%     watching/seeing as such (seeing = Aristotle's own example of
%     COMPLETE energeia, Metaph. IX.6 — a strict examiner would catch
%     the conflation).)

A world built for recreation never has to say why it should be entered.
Its player arrives already seeking something---an evening's diversion, a
story, a game worth playing---and the world reserves its arts of
solicitation for the encounters the player did not seek. A course world
begins from the opposite condition. The student arrives because the
course requires it, and so the solicitation that recruits them cannot
come from inside the world at all. It comes earlier, and it comes in
language. From the outset we presented the environment as the place
where a procedure group could watch its assigned footage together and
talk while watching---group analysis conducted inside a shared clinical
space, a way of working the course offered nowhere else.\footnote{The
platform's published description names the affordance directly: the
environment ``supported real-time voice chat and screen sharing,
enabling collaborative team viewing'' for groups meeting inside the
virtual environment (``Scalable Virtual Reality'' 4).} That
announcement is the deployment's act of solicitation. It names the
reason to enter this world rather than the YouTube platform that
carries the same footage. What remains for the rendered world itself is
the quieter half of the work: not arresting attention but keeping it,
across the weeks of a term, in a world the student was assigned rather
than drawn to.

Judged by what it set out to sustain, that quieter campaign succeeded.
Three of every five students wrote something affirming the footage's
interest, the largest positive register in the collected data;
presence language came from those who entered the world; attention,
once gathered, was held. The deepest complaint in the data is not that
the world failed to draw anyone. The campaign goes silent at exactly one
point. Interest that has been brought to commitment wants something to
complete itself in, and completing it would require the world to stop
soliciting and begin tasking---to station an \textit{orekton}, a
determinate object for desire and action: something to find, perform,
or submit from inside the hospital itself. The world contains no such
station. There is no procedure to attempt in its rooms, no
needs-finding to capture where the footage plays, no deliverable the
world can receive. In the chain's terms, the deployed world reliably
carries a student to committed interest at A\textsubscript{3} and
offers no crossing to A\textsubscript{4}. The hail itself made this
silence hard to see. What we announced was a place for the group's
analysis, not an analysis the world would conduct; the task always
lived on the course's side, in deliverables the world never housed. A
pitch can succeed---students entered---while the world it recruited for
withholds the one thing that would complete the interest it sustains. A
world that can only be watched has a boredom half-life.

None of this makes boredom a property of the footage. The same virtual
hospital drew the survey's fullest presence praise and its flattest
comparison, and its sharpest witness is a single student---the survey's
one confirmed co-watcher---who wrote in one answer that the platform
``does put you into a VR clinical environment'' where ``you can watch
with others'' (W26-R037/Q10), and in the next that it ``is not
significantly more immersive then a YouTube video'' (W26-R037/Q9).
Nothing about the hospital changed between those two answers; the
relation did. What holds for the footage holds for the route as well. A
student who tires of a long surgical recording will tire of it in the
virtual hospital and in a browser tab alike; I take the two routes to
differ in how soon that tiring arrives and in how far it goes, not in
whether it comes. Boredom, then, cannot be what separates the two
routes. Whatever it diagnoses, it diagnoses the watching itself---the
situation of being given something to see and nothing to do, which the
deployment and the YouTube platform stage alike. Saying what that
boredom is, rather than merely that it occurred, requires distinctions
of structure rather than degree. In the dissertation's account, boredom
takes three forms; the forms differ in what they attach to, and they
differ just as sharply in what evidence can show them.

Heidegger's lecture course of 1929--30, published as \textit{The
Fundamental Concepts of Metaphysics}, gives boredom one of the most
sustained, philosophical, and structural analyses it has received, and
the analysis begins from a definition small enough to carry: what
bores is ``that which holds us in limbo and yet leaves us empty''
(\textit{Fundamental Concepts} 87). Boredom, on this account, is
always two things at once. Something holds us: we cannot simply leave,
because the situation keeps us in it. At the same time, what holds us
offers nothing, so that we are left empty while being held. Against this
double grip we deploy a counter-move Heidegger calls passing the time
(\textit{Zeitvertreib}): occupations taken up not for their own sake
but to keep the boredom off. The three forms he distinguishes are
three depths to which this machinery can reach. Setting them out
requires one term of method. Call a channel any independent way an
experience shows itself in evidence: what a person says or writes
about their experience is one channel; what their visible behavior
shows is a second; what an instrument reads off the body is a third.
The forms of boredom differ in how many channels it takes to see them,
and for the collected data that difference matters as much as the
forms themselves.

In the first and shallowest form, we become bored \textit{by}
something. A definite thing is the culprit---a book, a lecture, a
wait---and we could name it if asked, because the boredom stays
attached to its object. Heidegger's scene is a missed connection at
``the tasteless station of some lonely minor railway'': ``It is four
hours until the next train arrives'' (93). The wait holds us---we
cannot leave without losing the train---and the station offers
nothing, and so we fight, in the open: we read the timetable, count
trees, draw figures in the sand, and ``catch ourselves looking at our
watch yet again'' (93). The repeated glance at the watch is what
Heidegger calls a ``helpless gesture'' (97): it shows that the passing
of time is failing against the wait. Everything about this form is
visible, and that is its evidential mark. Because the culprit can be
named and the fight can be seen, a single channel suffices to attest
it. A student who writes ``The VR is a cool idea but the quality is so
low I don't gain from the atmosphere and I dont feel immersed''
(W25-R014/Q9) has named the culprit inside the complaint itself; no
second instrument is needed to confirm what the sentence already
carries. The collected data holds this form's whole anatomy in
that one channel. The culprits are named: the footage's quality, its
length, the lag. The occasions are named: the crash at startup, the
buffering pause, the room that could not be found. The counter-move is
named too, because in this deployment the readiest way
to pass the time spent waiting on a struggling world was to leave it
for the same footage elsewhere: ``the youtube format ended up being
simpler to use and watch'' (W25-R017/Q9). Culprit, fight, and exit,
all written down in so many words.

The second form is stranger: we are bored \textit{with} something, and
no culprit can be found. Heidegger's scene is a dinner party. The food
is good, the talk is lively, nothing drags; asked at the table, we
would say, sincerely, that we are having a fine evening, and looking
back on it we can confirm that ``[t]here is nothing at all to be found
that might have been boring'' (110). Only afterward, at home, does the
recognition arrive: ``I was bored after all this evening'' (110).
Nothing was wrong with the evening, and that is the finding itself.
Where the first form's emptiness comes from a thing that refuses us,
here ``what is boring does not come from outside: it arises from out
of Dasein itself'' (128): the emptiness forms in us, underneath a
surface that remains genuinely, pleasantly occupied. The form has a
scene any student would recognize: an evening of episode after
episode, each one actually watched, none of them bad, which leaves one
hollowed rather than restored---nothing to blame, and yet the whole of
it empty. This structure---an occupied surface over a hollow depth---
is what makes the second form hard to evidence. The surface, asked,
would deny boredom and mean it; the depth cannot be asked. One channel
can therefore never show this form: a written answer reports the
surface, and the surface really is occupied. It becomes visible only
where two channels can disagree about one and the same experience, as
in the laboratory arm of this study, where a brain measure can read
boring-like on the very footage a gaze measure reads as engaged. A
course survey holds one channel. It can catch, at most, the form's
echo---praise of an occupation that works, standing over an emptiness
the student cannot name, the shape of the survey's two culprit-less
responses\footnote{``The virtual reality environment, while it allows
for more immersion, wasn't the highlight of the course, and could be
amended for something else'' (W26-R016/Q6); ``It is a cool concept,
but I feel like at the moment it is not doing much to help regarding
this class or the assignment'' (S25-R009/Q10). The insight the pair
provides is structural. Each response commends the platform---the
first even grants it the immersion it promises---and each reports an
emptiness without naming anything broken. The hollowness attaches to
no defect but to the whole (``wasn't the highlight,'' ``not doing much
to help''), and in the second response it attaches to what the
environment is \textit{for}: its help toward the class and the
assignment. That is the second form's silhouette in survey prose: the
occupied surface speaks well of the occupation while the emptiness
shows through only as a want the student cannot locate. With one
channel it remains a silhouette; nothing in either response can
confirm what stands beneath it.}---and that is why the deployment's
divergence register was read as echoes, never as detections.

The third form is the rarest and deepest, and it strips away even the
person who is bored. Its name is impersonal: ``it is boring for one''
(135)---not for me, with my roles and plans and reasons, but for one,
for anyone at all. No thing is the culprit here and no occasion is
required, because what goes flat is not this or that offering but
beings as a whole: everything at once, oneself included, becomes
indifferent. Heidegger's example is chosen for its very lack of
incident: it is boring for one ``to walk through the streets of a
large city on a Sunday afternoon'' (135). The example works because
nothing in it is boring. The shops are shut, but one wanted nothing
from them; there is no wait, no missed train, no evening to see
through; the afternoon lacks nothing in particular---and precisely
there, with nothing particular refusing us, everything refuses at
once. A student's version is near to hand: a Sunday on which the
assignment, the shows, the group chat, the whole of one's own room
stand available, and none of it---not even oneself---holds any pull,
and looking for something to do feels pointless before the looking
starts. That last detail is the form's signature. In the first form we
fight the boredom; in the second we fail to notice it; here we do not
even attempt the counter-move, because ``all passing the time is
powerless against this boredom'' (135)---one is instead ``compelled to
listen'' to what it discloses (136). Its evidence is accordingly
stillness and time: a comportment that has stopped searching for
occupation, and a while that ``becomes long'' (152), felt duration
stretching far past the clock. A retrospective course survey can
capture neither, and this analysis claims no instance of the form.
Where anything of its shape appears in this study's evidence---in the
laboratory's stillest episodes and dilated felt-time reports---it is
read under an announced caution, as analog rather than instance.

Read with this apparatus, the collected data sorts almost
entirely into the first form---honestly so, since a survey is a
single channel and the first form is the one a single channel can
hold. What the apparatus adds is the structure beneath the sorted
complaints. Heidegger names the two moments that constitute every
form: being left empty (\textit{Leergelassenheit}) and being held in
limbo (\textit{Hingehaltenheit}). His most exact observation about the
first is that emptiness is not absence: ``To leave empty means to be
something at hand that offers nothing'' (103); things ``leave us empty
precisely because they are at hand'' (102). The virtual hospital is at
hand in every sense that matters---rendered, explorable, stocked with
footage, praised in the survey's fullest presence language---and
precisely as such it leaves empty, because what it puts at hand offers
everything to see and nothing to complete. The station's refusal even
has an address: ``The station at hand refuses itself to us as a
station and leaves us empty because the train that belongs to it has
not yet arrived'' (103). The virtual hospital refuses itself as a
clinic in exactly this way. A clinic is a place where procedures are
performed and needs are found; this one holds only their images, and
its train never arrives---no task pulls in to make the place what its
rooms announce. One student stated the resulting emptiness with
precision: the procedures are ``many hours long, and I'm not sure what
exactly to be looking for during that time'' (S25-R035/Q6)---at hand,
and offering nothing determinate to look \textit{for}. The second
moment completes the situation. At the station we are held by an
interval with a terminus: four hours, then the train. The deployment's
students are held by mandate, with no terminus inside the world
itself---assigned footage, required hours, a term's duration, and, in
the small, the buffering pauses that hold a viewer between intention
and image. Held before an offer that gives nothing to complete: the
two moments meet, and they meet at one point only, the world's last
link.

The diagnosis follows, and it is a diagnosis of a situation, not of a
platform. Boredom, as the collected data presents it, marks a
mis-attunement among three terms: a learner who arrives under mandate,
with no leisure motive for the world to ride on; an ecology that does
its sustaining work, as the data shows it doing; and a task that has
no address inside the world. The interest the world sustains is not a
contemplative looking that completes itself in every moment of its
exercise. It is interest underway toward doing---the students'
requests for assignments, for
interaction, for something to find make its direction explicit---and
motion of that kind is, in Aristotle's terms, \textit{energeia
atelēs}: incomplete actuality, actuality whose end lies beyond it
(\textit{Physics} III.1, 201a10). A world that sustains the motion
while housing no terminus does not merely disappoint the interest; it
holds it, unfinished, for as long as the mandate runs. The
deployment's shortfall, then, was never bandwidth---presence survived
the monitor---but completion. This is the precise sense in which the
deployed world's emptiness is regional. Nothing upstream is withheld:
the world is entered, rendered, watched, even dwelt in; the
being-left-empty is stationed at the chain's last link, where
soliciting would have to become tasking and the world falls silent.
The boredom the data attests is therefore not a verdict on the
footage, and not a difference between the two routes; it is the felt
shape of the watching situation both routes stage, and, unlike a crash
or a buffer, it names nothing a repair could answer. The arc of the
data shows as much: the frictions fell term by term, and what rose in
their place was the question of what the world is for. An emptiness
stationed where action should begin asks for one thing only---a
for-the-sake-of-which, a reason this world, rather than a video, is
where the work of the course happens.

% --- M4.2 Veri(dis)similitude's second face (plan approved w/ P1 rulings
%     2026-07-27: (1) opening = plain restatement of the revised
%     veri(dis)similitude definition, NO stylistic metalanguage ("the
%     deployment's twist" etc. banned); (2) hypothesis notation (H1/H3)
%     NEVER in rendered text (extends the module-label rule); (3) MA-thesis
%     poles RESTATED not quoted (decision point resolved by ruling 1).
%     4-¶ architecture: ¶1 condition restated + turned inward · ¶2 cohort
%     re-derives it + the un-redeemed answer · ¶3 materials audit
%     (commodity features → YouTube; place strongest form, insufficient
%     alone) · ¶4 the material video cannot approximate (hands M4.3 its
%     premise; de-facto scoping stays in M4.3). ¶¶1–2 DRAFTED 2026-07-27 —
%     in review. Verbatims reused from verified ledger: W26-R033/Q9;
%     W26-R022/Q10 ("the least gimmicky" + meet-up proposal, same
%     response per M3.6). Part II two-pole sentence ("too familiar cannot
%     move / too strange cannot be entered") reused as internal formula,
%     no cite. Nine-campus facts = ¶3 (M0.2 + author-confirmed memory).
%     ¶¶1–2 v2 RULINGS APPLIED 2026-07-27: ¶1 close EXPANDED per user
%     (subtract-the-footage comparison of the two remainders + the
%     film/theater demonstrative example); ¶2 opening = user formula
%     ("As the collected data demonstrates"); "RECORD" BAN RE-AFFIRMED
%     (standing framing rule) → GLOBAL SWEEP executed across M4.1
%     approved text + M4.2 (16 prose instances → collected data/data/
%     survey/students' language; wording-only touches to approved ¶¶,
%     And-precedent). NB the flagged "remove first sentence" items were
%     chat-presentation labels, never in this file; future batch
%     presentations carry bare ¶-numbers only.
%     ¶¶1–2 APPROVED 2026-07-27 w/ one ordered change APPLIED: the term
%     veri(dis)similitude stated DIRECTLY (named at the definition, ¶1
%     opener, verisimilar/dissimilar poles matched to the term; and at
%     the turn "Veri(dis)similitude accordingly takes a second and
%     sharper form"). ¶¶3–4 DRAFTED 2026-07-27 — in review (materials
%     audit + the unapproximable material; W25-R014/Q6 quality-trade
%     quote char-exact vs open_text.csv, inner ""immersion"" → MLA
%     ¶4 v2 APPROVED 2026-07-27 — **M4.2 COMPLETE (4 ¶¶, all approved)**.
%     single quotes, extends the W25-R014 thread [Q9 in ¶5 of M4.1];
%     ASEE fallback span "with interactive 360-degree panning" verified
%     char-exact via pdftotext at PDF p. 4, cite ("Scalable Virtual
%     Reality" 4); nine-campus facts per M0.2 + author-confirmed memory;
%     hands M4.3 its premise, de-facto scoping reserved to M4.3.
%     ¶3 APPROVED 2026-07-27. ¶4 v2 per ruling 2026-07-27: opening two
%     sentences REWRITTEN slow+concrete ("rival route"/"rides either
%     route" compression replaced by worked examples — course-side
%     assignment [worksheet/quiz/unmet-needs prompt] works under either
%     delivery because administered outside the footage; in-world
%     needs-task = unbuilt possibility; browser panning = drag-the-view,
%     lacks corridor/doorway/position); "approximation" → "imitation"
%     diction unified (one-word touch to the approved close). IN REVIEW.)

A rendered world, if it is to do anything to the person who enters it,
must satisfy a double condition: \textit{veri(dis)similitude}. It must
be verisimilar enough---familiar enough in its look, its workings, its
ways of being moved through---that a person can enter it and dwell in
it at all; and it must be dissimilar enough---different enough from
what everyday life already offers---that the entering changes
something. A world too familiar cannot move; a world too strange
cannot be entered. For the deployed virtual hospital, the second half
of this condition is measured against an unusual standard. Every
procedure the hospital holds also plays on the YouTube platform, one
click away, free of installers, sign-ins, and crashes. A student
deciding whether to enter the world is therefore not weighing it
against everyday reality; they are weighing it against the video of
itself. Veri(dis)similitude accordingly takes a second and sharper form.
Ordinarily a rendered world earns its dissimilarity against everyday
life, out of possibilities everyday life does not hold. The deployed hospital
cannot earn its difference there, because its substance---the
procedures---exists identically outside it. Subtract the footage from
both routes and compare what remains. On the YouTube platform the
remainder is a familiar player and its conveniences: the progress bar,
the subtitles, the speed control. In the virtual hospital the
remainder is everything the footage plays inside of: the walk down the
corridor, the room found and entered, the screen on its wall, and the
possibility of another student standing at the same screen. The world
justifies its existence with that remainder or not at all. The demand
is not unique to virtual worlds. A film that can be streamed at home
identically must give a viewer a reason to go to the theater, and the
theater's answer has never been the film, which is the same in both
places; the answer is the hall, the occasion, and the others in the
seats. So with the deployed world: to justify its existence it must be
different enough from the video of itself, and the difference can be
made only of what surrounds the footage---never of the footage, which
both routes share.

As the collected data demonstrates, the cohort arrived at this same
condition from inside it. As the platform stabilized and the
complaints about function fell away, the talk that rose in their place
questioned existence rather than function: ``I think the youtube
videos were enough to fully grasp the experience'' (W26-R033/Q9). A
question of that kind, put to a world whose footage plays elsewhere,
is the difference-condition stated in a student's words. The question
was not unanswered. It had been answered before the first term began,
in the announcement that recruited the class: the difference was
company---a procedure group watching together, inside the world,
talking while it watched. What the data then shows is that
difference going largely unconvened. Students watched alone, found
what a world minus its difference amounts to, and asked why it
existed. The survey's most skeptical voice---the student who judged
traditional learning ``the least gimmicky'' way---proposed, in the
same response, ``perhaps the ability to meet up with groups in a
virtual reality setting'' (W26-R022/Q10): the difference itself,
proposed back to the course that had announced it. The question the
deployment ends on is an answer it began with, delivered and never
redeemed.

What could the difference be made of? Not the commodity features of
delivery, which the students' own comparisons settle. The YouTube
platform streams the footage at higher quality for the same bandwidth,
carries subtitles, and changes playback speed at will; one student
weighed the trade directly, having ``gained more from the youtube
videos because I could watch in much higher quality which made up for
the slight reduction in `immersion'{}'' (W25-R014/Q6). On every
feature a video player can have, the player built for videos wins.
Three materials remain that no player has: a task, a place, and
company. A task the deployed world did not contain; that absence is
the diagnosis already made. Place it did contain, and the place
deserves its full weight. This class has never had a room. Its
students enroll across nine university campuses whose calendars do not
even align; there is no lecture hall, no laboratory section, no
building, no campus they share. The virtual hospital is the only room
the class has ever had---not a supplement to a classroom but the sole
site at which this cohort could be co-present in any sense at all. Yet
the collected data shows place alone not sufficing. Presence was had
on the desktop route and attested plainly, and the flight to the
YouTube platform happened anyway; a place entered alone, with nothing
to do in it, reduces in use toward the very thing it competes
with---a scene to survey, reached the harder way.

Two of the three materials are not wholly beyond the YouTube
platform's reach. Consider the task first. The course can attach an
assignment to the videos wherever they play: a worksheet completed
while watching, a quiz on the procedure, a prompt to list the unmet
needs one observed. An assignment of that kind works as well under a
YouTube link as inside the virtual hospital, because the course
administers it and the student completes it outside the footage---no
part of it needs the world. A world could go further: it could make
the finding of needs something done in its rooms, captured where the
footage plays. The deployment built no such task, and until one is
built, the task-material does not separate the two ways of watching.
The place can be partly imitated as well. The footage itself is
360-degree film, and the deployment's own fallback delivered the same
videos ``with interactive 360-degree panning'' (``Scalable Virtual
Reality'' 4): a student in a browser can drag the view around the
operating room and look wherever they choose. What the browser cannot
give is everything around that view---the corridor, the doorway, a
position in the room---so the pannable video is a thin, surveyable
shadow of the place; but it is a shadow of it. The meeting admits no
imitation at all. Watching together in a shared rendered room,
talking while the procedure plays, is not a feature a video player
lacks in degree; it is a kind of thing no video player has. For most
pairs of students in this class there is no library to meet at
instead, no dorm lounge, no lab bench; the alternative to meeting in
the virtual hospital is not meeting somewhere else---it is not meeting
at all. If the world's difference from the video of itself is to be
made of anything, it must be made of the meeting.

% --- M4.3 The WE as the desktop's ontological differentiator (★ SPOTLIGHT.
%     Plan approved 2026-07-27 incl. the ONE-SENTENCE marked-speculative
%     WE–boredom acknowledgement in ¶5. 5-¶ architecture: ¶1 dual-being +
%     toward-whom ground · ¶2 route audit + de-facto scope · ¶3 being-with
%     mode + honesty clause · ¶4 deferred theorizing + floors · ¶5 climax.
%     ¶¶1–2 DRAFTED 2026-07-27 — in review.
%     Batch-1 verifications: "multiplayer functionality, voice chat, and
%     customizable avatars" = AUTHORS' OWN prose ("our learning
%     environments"), pdftotext-pinned PDF p. 10 = journal p. 390 →
%     first-person rendering, cite ("Phenomenological Evaluation" 390);
%     p. 394 multiplayer hit = appendix instrument text, NOT used.
%     Toward-whom = Part I canonical (toward-which/toward-whom intentional
%     structure; thymotic action "always in relation to another"; Mitsein
%     slot) — RESTATED without cite per no-re-derivation + ambiguous
%     Bekker ranges across Part I files (1378a31–b5 vs a30–b9; NB
%     FourTypes Type-2 sentence carries a PROVISIONAL marker — nothing
%     quoted from it). Burke NAMED for consubstantiation, no quote —
%     ADD A Rhetoric of Motives to Works Cited at assembly.
%     M1.6 inheritance (cite forms per its locked header): Biocca et al.;
%     Terry and Doolittle — argued use lands in ¶3.
%     ¶¶3–5 APPROVED 2026-07-27 — **M4.3 COMPLETE (5 ¶¶, all approved)**.
%     Batch-2 verifications:
%     Biocca "the sense of being with another" (456) + Terry & Doolittle
%     relocation (121, paraphrase-cite) + Kreijns Pitfall 1 (336,
%     paraphrase-cite; full deployment reserved for M5.2) = M1 bank
%     PDF-verified reuses, cite forms per M1.6; W25-R038/Q10 "when I did
%     log in, I was usually alone in the platform...connected with my
%     peers" char-exact vs open_text.csv this session; W26-R037/Q10
%     fragment = ledger reuse (continuity). Floors per M3.5 locked spec
%     (1/11/4/13, awareness 12) — counts only, NO rates. ¶5 carries the
%     approved ONE-SENTENCE marked-speculative WE–boredom acknowledgement
%     + answers M3.4's open question (exclusive binding = the others) +
%     un-redeemed-hail close ("temporal economy" = outline-approved
%     phrasing).)

Rhetorical incorporation is one in substance and dual in being. A
world becomes a student's own in two ways at once: as a place they
dwell in---world-disclosedness, the register every presence finding so
far has measured---and as a place that is theirs together with
others---co-disclosedness, the same world had as ours. The two are not
a feature and its supplement; they are two sides of one
phenomenon, and the chain itself grounds the second. Desire's thymotic
species is irreducibly social: one is angry at someone, stands on
honour before someone; the emotions that concretize judgment into act
are articulated toward or before determinate others---anger toward the
slighter, shame before the audience that matters. An act, when it
finally comes, is rarely aimed at no one; the chain's last link
carries a toward-whom as part of its structure. What Burke names
consubstantiation---persons made, through acting together, of one
substance---names the same ground from rhetoric's side: to act with
another is to be with them in a way no shared spectacle alone
provides. A world that supports action will support it toward and with
others, or it will support less than action's full structure.

The three routes hold this dual structure very differently. The
published headset version maximized world-disclosedness: the encircling
display, the fullest supply of the operating room a medium could give.
Co-disclosedness it did not carry. Students entered alone, and we said so in
print, naming among the version's limits ``the sheer lack of features
found in traditional games like multiplayer functionality, voice chat,
and customizable avatars'' (``Phenomenological Evaluation'' 390). The
YouTube platform possesses neither in kind: it is a player, not a
world; there is nothing in it to be in, together or otherwise. The
desktop route inverts the headset's profile. Its world-disclosedness
rests on the thinnest supply: the monitor gives less, and \textit{phantasia} furnishes
more of the world than on any other route. But as the program stands
deployed, it is the only route on which co-disclosedness is available
at all. The scope of that claim is historical, not metaphysical:
proximity voice could be built into a headset version---nothing in the
hardware forbids it---but development moved to the desktop build, and
the meeting was built there, only there. The affordance is not a
consolation native to screens; it is a developed asset this route
alone carries. The route that supplies the least world is the route
that can share it.

What the meeting changes is not the platform's feature list but the
being of the world. A room entered alone is a viewing station; the
same room entered with another student is ours before anything is said
in it---a place jointly attended, where the judgment forming over the
footage can be spoken, answered, and made common. The research on
learning at a distance has a name for this register of experience and
a hard-won result about it. Social presence---defined at its core as
``the sense of being with another'' (Biocca et al.\ 456)---is not a
property a medium contains but a state persons achieve, something a
course can activate or let lie (Terry and Doolittle 121). Those
findings were made about text forums and video calls, not about
rendered worlds; the two research traditions descend from a common
origin and have scarcely cited one another since. Joining
them---reading a rendered world's meeting through being-with---is this
dissertation's own synthesis, offered as such. What it yields, if it
holds, is exact: being-with is a second presence, of its own kind,
that the desktop hospital can host---and the first kind, presence in a
place, cannot substitute for it.

A world, though, can only hold a meeting open; it cannot convene one.
Footage is delivered to whoever presses play, and a place is there for
whoever walks in, but being-with requires that others arrive. Arrival
is exactly what the deployment's shape made unlikely. Convening a
group asks what solo watching never asks: schedules aligned across
campuses that share no calendar, an invitation risked on classmates
never met, an hour kept. The course was built so that no student ever
needs to share an hour with another---that is what systemwide
asynchrony means---and every convenience the YouTube comparison
rewards pulls down the same solo path. The collaborative-learning
literature warns of precisely this: interaction cannot be assumed
merely because an environment makes it possible (Kreijns et al.\ 336).
The collected data shows the warning borne out. In 873 responses there
is one confirmed convening---the student who wrote that ``you can
watch with others'' (W26-R037/Q10). Eleven students wrote of being
alone in the world; ``when I did log in, I was usually alone in the
platform,'' one of them explained, ``so it didn't really make me feel
more connected with my peers'' (W25-R038/Q10). Four asked the course
to require the meeting; thirteen asked for the capability, twelve of
them in wording implying it did not exist. The affordance that answers
the world's justification-question was unknown to part of the class
that wanted it, unconvened by nearly all of it, and lived by one voice
in three terms. The world held the WE open; the course's temporal
shape let it lapse.

The deployment's deepest loss now has its exact name. It is not
presence---presence the desktop world had, measured and attested. It
is the WE left unconvened: a co-disclosedness the ecology held open
for three terms and the temporal economy let lapse. The arc that
began this analysis closes on itself. The difference we announced at
the course's outset was the meeting; the audit of materials arrived
at the meeting; the survey's most skeptical voice, unprompted,
reinvented the meeting. The announcement was right, and what it named
went largely unlived. The question the presence analysis left
open---what could bind a world to a student exclusively, when its
footage plays anywhere---has the same answer: its others. The virtual
hospital is irreplaceable exactly insofar as it is shared, and
replaceable exactly insofar as it is not. Whether a convened WE would
also blunt the boredom this analysis diagnosed---company giving
sustained interest a completion of another kind---the collected data
cannot say; the suggestion is speculative, and it is offered as such.
What the data can say is narrower and firmer: the one thing this
world alone could be was the one thing it was almost never used as.

% --- M4.4 What the desktop case gives back to the frame (SOFT-BLOCK
%     WAIVED by user 2026-07-27 — the Part II closing differential is not
%     yet written. ***REVISION-PASS REQUIRED: revisit this sub-module
%     once the Part II closing differential exists; harmonize the
%     empirical register here against the differential's final claims
%     (division of labor AUTHOR-CONFIRMED 2026-07-16: Part II = media-
%     type/interaction-mechanics/design differential incl. phantasia's
%     near-leveling; Part III = diagnosis of the experienced difference).***
%     2 ¶¶ per outline: frame-return (matched numbers 70%/70.7% INHERITED
%     from M3.4, not re-derived; phantasia-leveling echo AUTHOR-APPROVED
%     2026-07-16; "Route sets the load; the task sets the dwelling" =
%     cross-case §2(vii) motto, first rendered appearance) + one
%     gesture-scale ¶ on what only a headset deployment could test
%     (gesture only per standing rule). Terminology lock held ("phantasia
%     furnishes more" — no image-making/phantasia-load). ¶¶1–2 APPROVED
%     2026-07-28 — **M4.4 COMPLETE. M4 MODULE COMPLETE (M4.1 9¶¶ ·
%     M4.2 4¶¶ · M4.3 5¶¶ · M4.4 2¶¶ — all approved). M4.4 revision-pass
%     pending Part II closing differential (see marker above).**)

Two findings from this deployment reach beyond it. The first is the
pair of matched numbers: the pilot's 70\% under headsets, the
deployment's 70.7\% before ordinary monitors. The supply could hardly
differ more---an encircling display against a window on a desk---yet
the felt result barely moved. Parity of that kind is what the leveling
power of \textit{phantasia} looks like in data: where the medium
supplies less, \textit{phantasia} furnishes more, and the felt world
arrives nearly whole either way. The second finding is what no route
could level: the act that none of them contained, and the meeting that
only one of them carried. Route sets the load; the task sets the
dwelling. A medium decides how much of the world \textit{phantasia}
must supply; whether anyone comes to dwell in that world---and with
whom---is decided at the last link, by what there is to do there and
who there is to do it with. At the scale of 241 students, the
deployment attests that this is where worlds are won or lost: at the
crossing from A\textsubscript{3} to A\textsubscript{4}, where
committed interest either finds its act or watches instead.

What this case could not test, a headset deployment at the same scale
could: whether the parity of presence holds when the fullest supply
meets a task actually built into the world, and whether convening
comes easier to students embodied together in an encircling world than
to students meeting through windows on desks. Those are questions for
a version of the course that does not yet exist, and they remain
questions here.
